\documentclass[12pt, a4paper, oneside]{book}
\usepackage[hidelinks]{hyperref}


\usepackage[english,slovak]{babel}	
%\usepackage{epsfig}
%\usepackage{epstopdf}
\usepackage{amsfonts}

%\usepackage[chapter]{algorithm}

\usepackage{algorithm}
%\usepackage{algorithmic}
\usepackage{algpseudocode}

\usepackage{listings}
%\usepackage{amsmath}
%\usepackage{amssymb}
\usepackage{graphicx}
\usepackage{multirow}
\usepackage{color}
\usepackage{url}
\usepackage[utf8]{inputenc}
\usepackage[T1]{fontenc}
\usepackage{setspace}
%\usepackage{tabularx}
\usepackage{textcomp}
\usepackage{caption}
\usepackage{natbib}

%\usepackage{amsmath}

\hyphenpenalty=10000 
\bibliographystyle{abbrv}

\setstretch{1.5}
%\renewcommand\baselinestretch{1.5} % riadkovanie jeden a pol

% pekne pokope definujeme potrebne udaje
\newcommand\mftitle{Abduction solver based on  highly   efficient C++ DL reasoner}
\newcommand\mfthesistype{Master thesis}
\newcommand\mfauthor{Bc. Drahomír Mrózek}
\newcommand\mfadvisor{RNDr. Martin Homola, PhD.}
\newcommand\mfconsultant{Mgr. Júlia Pukancová}
\newcommand\mfplacedate{Bratislava, 2017}
\newcommand\mfuniversity{Comenius university in bratislava}
\newcommand\mffaculty{Faculty of Mathematics, Physics and Informatics}
\newcommand{\sub}[1]{$_{\text{#1}}$}
\newcommand{\reference}[1]{č.~\ref{#1}}
\newcommand{\imageHeight}{150px}

\ifx\pdfoutput\undefined\relax\else\pdfinfo{ /Title (\mftitle) /Author (\mfauthor) /Creator (PDFLaTeX) } \fi
\newtheorem{mytheo}{theorem}[section]

%\theoremstyle{definition}


\newtheorem{mydef}{Definition}[section]
\newtheorem{myex}{Example}[section]
\begin{document}

\frontmatter

\thispagestyle{empty}

\noindent
\begin{minipage}{\textwidth}
\begin{center}
\textbf{\mfuniversity \\
\mffaculty}
\end{center}
\end{minipage}

\vfill
\begin{figure}[!hbt]
	\begin{center}

		\label{img:logo}
	\end{center}
\end{figure}
\begin{center}

		{\textbf{\Large\MakeUppercase{\mftitle}}}\\
		\smallskip
		{\mfthesistype}
	
\end{center}
\vfill
2018 \hfill
\mfauthor
\eject 
% koniec obalu

\thispagestyle{empty}

\noindent
\begin{minipage}{\textwidth}
\begin{center}
\textbf{\mfuniversity \\
\mffaculty}
\end{center}
\end{minipage}

\vfill
\begin{figure}[!hbt]
\begin{center}
\includegraphics{images/logo_fmph_dark}
\label{img:logo_dark}
\end{center}
\end{figure}
\begin{center}
{\textbf{\Large\MakeUppercase{\mftitle}}}\\
		\smallskip
		{\mfthesistype}
\end{center}
\vfill
\begin{tabular}{l l}
%Registration number: & 40a99bd8-3cb6-4534-9330-c7fd9b5e5ca4 \\
Study programme: & Applied informatics\\
Field of study: & 2511 Aplikovaná informatika\\
Study department: & Department of Applied Informatics\\
Advisor: &  \mfadvisor \\
Consultant: & \mfconsultant
\end{tabular}
\vfill
\noindent
\mfplacedate \hfill
\mfauthor
\eject 
% koniec titulneho listu

%\thispagestyle{empty}
%\includegraphics[width=\textwidth]{images/zadanie}
%\vfill
%\eject
% koniec zadania

\thispagestyle{empty}


\begin{figure}[H]
\begin{center}
todo:zadanie
\label{img:zadanie}
\end{center}
\end{figure}

{~}\vspace{12cm}

\noindent
\begin{minipage}{0.25\textwidth}~\end{minipage}
\begin{minipage}{0.75\textwidth}
Čestne prehlasujem, že túto diplomovú prácu som vypracoval samostatne len s použitím uvedenej literatúry a za pomoci konzultácií u môjho školiteľa.
\newline \newline
\end{minipage}
\vfill
~ \hfill {\hbox to 6cm{\dotfill}} \\
\mfplacedate \hfill \mfauthor
\vfill\eject 
% koniec prehlasenia

\chapter*{Acknowledgments}\label{chap:thank_you}
TODO
\vfill\eject 
% koniec podakovania
\chapter*{Abstract}\label{chap:abstract_en}
This work will start from an existing solution of an abduction solver based on Pellet reasoner for description logics, implemented in Java. The implementation in C++ opens different ways for improvement and effectivization, starting with the utilization of a more effective C++ inference system for description logics.
~\\
Keywords: abduction, description logic
\vfill\eject 
\chapter*{Abstrakt}\label{chap:abstract_sk}
Práca bude vychádzať z existujúceho riešenia abduktívneho systému založeného na reasoneri pre deskripčné logiky Pellet implementovaného v Jave. Pri implementácii v C++ sa otvárajú možnosti na zlepšenie a zefektívnenie existujúceho návrhu abduktívneho solvera, počnúc využitím efektívnejšieho C++ inferenčného systému pre deskripčné logiky.

~\\
Kľúčové slová: abdukcia, deskripčná logika
\vfill\eject 




% koniec abstraktov
\renewcommand{\contentsname}{Contents}
\tableofcontents

\mainmatter

\renewcommand{\chaptername}{Chapter}

% treba este prejst dokument ci je kod spravne formatovany
\input 01intro.tex
\input 02description_logic
\input 03abduction
\input 04previous_solutions
\input 05our_solution
\input 06implementation
\input 07results
\input 08conclusion
\backmatter

\nocite{*}
%\bibliographystyle{abbrv}
\bibliography{references}

\listoffigures

\end{document}